\documentclass[11pt,a4paper]{article}

\usepackage{kotex}
\usepackage[margin=22mm]{geometry}
\usepackage{amsmath,amssymb}
\usepackage{booktabs}
\usepackage{longtable}
\usepackage{array}
\usepackage{enumitem}
\usepackage{setspace}
\usepackage{xcolor}

\setstretch{1.16}
\setlength{\parindent}{1em}
\setlength{\parskip}{0.45em}
\renewcommand{\arraystretch}{1.25}

\title{\textbf{5PL 미들마일 플랫폼 연구의 문제 설정 및 시장 시나리오 설계}}
\author{유재우}
\date{2026년 9월}

\begin{document}

\maketitle

\section{문서의 목적}

본 연구는 5PL 미들마일 플랫폼의 강건 선제계약 및 최적 경매주기를 모델링하며 핵심 연구질문은 다음과 같다.

\begin{quote}
\textbf{불확실한 수요와 운송망 충격 하에서 5PL 미들마일 플랫폼은 얼마나 많은 운송용량을 선제계약하고,
경매를 얼마나 자주 열어야 하는가?}
\end{quote}

본 연구는 가격결정 자체를 주된 연구대상으로 삼지 않는다.
매칭은 판매자, 소비자, 운송업체 및 품목의 공급ㆍ수요ㆍ운송용량을 바탕으로 수행하고,
가격청산은 매칭 결과 이후 MMTR--MVCG 형식의 사후 절차로 처리한다.
시장 시나리오 역시 운송업체의 지역 선호, 거리 선호, 특정 품목의 납기 우선순위, 전략적 가격인상과 같은 추가 행동을 도입하지 않는다.

\section{시장 시나리오 설계}

\subsection{시장논리에 기반한 매칭}

본 연구에서는 특정 소비자 또는 품목에 인위적인 우선순위를 부여하지 않는다.
예를 들어 긴급화물, 고가화물, 특정 납기 민감품목을 별도로 정의하여 우선 배정하는 구조는 사용하지 않는다.
두 품목이 존재하더라도 두 품목의 차이는 중량과 부피와 같은 \textbf{물리적 특성}에만 둔다.

또한 Load-to-truck ratio 사례와 같은 거리, 지역, 특정 노선의 선호도와 같은 공간적 가중치도 도입하지 않는다.
따라서 특정 지역을 운송업체가 선호하거나, 특정 판매자--소비자 조합에 지역적 이점이 존재한다는 가정은 배제한다.
\textcolor{red}{본 연구에서 관찰하고자 하는 현상은 ``어느 지역이 유리한가''가 아니라 ``주어진 공급, 수요, 운송용량 하에서 어떤 거래조합이 최적의 매칭을 형성하는가''이다.}

만약 경매주기 또는 강건성 수준을 변경했을 때 성과가 달라진다면,
그 원인을 지역 선호나 품목 우선순위가 아니라
\textbf{수급균형, 운송용량, 선제계약 수준 및 불확실성}에서 찾을 수 있다.

\begin{table}
\centering
\caption{본 실험에서 의도적으로 제외하는 요소}
\begin{tabular}{p{0.30\textwidth}p{0.60\textwidth}}
\toprule
\textbf{요소} & \textbf{처리 원칙} \\
\midrule
품목별 우선순위 & 사용하지 않음. 품목은 물리적 특성만 다르게 설정 \\
거리 및 지역 가중치, 선호도 & 사용하지 않음 \\
공급부족에 따른 가격인상 & Core MILP에 포함하지 않음 \\
가격결정 & 매칭 이후 MMTR--MVCG 방식으로 사후 계산 \\
\bottomrule
\end{tabular}
\end{table}

\section{기본 시장구조}

\subsection{시장 참여자 비율}

시장 참여자 수는 판매자, 소비자, 운송업체의 상대적 규모를 일정하게 통제하는 방식으로 설정한다.
초기 기준비율은 다음과 같다.

\[
|I|:|J|:|K| = 1.5:2:1.
\]

이를 작은 정수 규모로 구현하면 $3:4:2$가 되지만,
수치실험에서 지나치게 작은 시장이 되는 것을 방지하기 위해 기준 실험에서는 두 배 규모인

\[
\boxed{|I|=6,\qquad |J|=8,\qquad |K|=4}
\]

를 기본값으로 사용한다.
품목은 두 개만 사용하여

\[
\boxed{|L|=2}
\]

로 둔다.
이 정도의 품목 수만으로도 중량 제약과 부피 제약이 서로 다른 방식으로 작동하는
weigh-out 및 cube-out 상황을 관찰할 수 있으므로,
초기 실험에서 불필요하게 품목 수를 늘리지 않는다.

\begin{table}[h!]
\centering
\caption{기준 시장의 참여자 구성}
\begin{tabular}{ccc}
\toprule
\textbf{구분} & \textbf{기호} & \textbf{기준 규모} \\
\midrule
판매자 & $I$ & 6 \\
소비자 & $J$ & 8 \\
운송업체 & $K$ & 4 \\
품목 & $L$ & 2 \\
\bottomrule
\end{tabular}
\end{table}

\subsection{품목 구성}

품목 $A$와 $B$는 우선순위가 아닌 물리적 특성만 서로 다르게 설정한다.
예를 들어 다음과 같은 대칭적인 구조를 사용할 수 있다.

\begin{table}[h!]
\centering
\caption{두 품목의 물리적 특성 예시}
\begin{tabular}{cccc}
\toprule
\textbf{품목} & \textbf{해석} & \textbf{중량계수 $w_l$} & \textbf{부피계수 $v_l$} \\
\midrule
$A$ & 중량소모가 상대적으로 큰 품목 & 10 & 2 \\
$B$ & 부피소모가 상대적으로 큰 품목 & 2 & 10 \\
\bottomrule
\end{tabular}
\end{table}

이 설정의 목적은 품목 $A$ 또는 $B$를 중요품목으로 만들기 위한 것이 아니다.
동일한 운송업체의 제한된 중량 및 부피 capacity를 두 품목이 서로 다른 비율로 사용하도록 만들어, 다품종의 거래를 표현하기 위함이다.

\section{판매자와 품목의 관계 설정}

판매자와 품목의 관계는 세 가지 방식으로 고려하였다.

\begin{longtable}{p{0.10\textwidth}p{0.28\textwidth}p{0.26\textwidth}p{0.26\textwidth}}
\caption{판매자--품목 관계의 대안 비교} \\
\toprule
\textbf{대안} & \textbf{구조} & \textbf{장점} & \textbf{판단} \\
\midrule
\endfirsthead
\toprule
\textbf{대안} & \textbf{구조} & \textbf{장점} & \textbf{판단} \\
\midrule
\endhead

1 &
한 판매자가 여러 품목을 동시에 공급.
예: 판매자 1이 $A$ 10개, $B$ 20개 공급 &
가장 일반적인 multi-item 구조이며 기존 $\hat s_{ilt}$를 그대로 활용 가능 &
가능한 구조이나 모든 판매자가 모든 품목을 공급하면 시장이 지나치게 dense해질 수 있음 \\

2 &
판매자별 판매가능 품목을 구분.
예: 판매자 1은 $A$만 10개, 판매자 2는 $B$만 20개 공급 &
판매자 전문화가 명확하며, 판매하지 않는 품목은 $\hat s_{ilt}=0$으로 표현 가능.
두 품목은 공통 운송용량을 두고 경쟁 &
\textcolor{blue}{\textbf{기준 시나리오로 채택}} \\

3 &
판매자 1명을 화물 1단위로 해석.
예: 판매자 1명이 $A$ 1개만 공급 &
개별 주문 단위의 이산 거래처럼 보일 수 있음 &
판매자(agent)와 공급수량(quantity)을 혼동하고 시장규모가 불필요하게 증가하므로 사용하지 않음 \\

\bottomrule
\end{longtable}

기준 시나리오에서는 판매자 6명을 품목별로 전문화한다.
예를 들어 판매자 1--3은 품목 $A$를 공급하고,
판매자 4--6은 품목 $B$를 공급하도록 설정한다.
판매하지 않는 품목의 공급량은 0으로 두면 되므로 새로운 이진 파라미터를 추가할 필요가 없다.
반면 모든 소비자는 두 품목 모두에 대한 수요를 가질 수 있도록 설정한다.
동일 품목에 대해 하나의 판매자는 복수의 소비자에게 공급할 수 있고,
하나의 소비자도 복수의 판매자로부터 물량을 공급받을 수 있도록 사전에 판매자--소비자 쌍을 고정하지 않는다.
따라서 품목 $A$와 $B$ 각각에서 3명의 판매자와 8명의 소비자가 다대다로 거래할 수 있으며,
본 연구의 매칭은 품목별 
\textbf{N:M matching 구조}를 갖는다.

\begin{table}[h!]
\centering
\caption{기준 판매자 전문화 구조}
\begin{tabular}{ccc}
\toprule
\textbf{판매자} & \textbf{품목 $A$ 공급} & \textbf{품목 $B$ 공급} \\
\midrule
$i_1,i_2,i_3$ & 양의 공급량 & 0 \\
$i_4,i_5,i_6$ & 0 & 양의 공급량 \\
\bottomrule
\end{tabular}
\end{table}

이 구조에서도 품목별 시장이 완전히 독립되는 것은 아니다.
두 품목은 동일한 운송업체의 중량 및 부피 capacity를 공유하므로,
실행단계의 매칭에서는 제한된 운송자원을 두고 간접적으로 경쟁하게 된다.
또한 판매자 또는 소비자별 거래상대의 수를 1개로 제한하는 별도의 제약을 두지 않음으로써,
실제 최적화 단계에서도 N:M 구조가 유지되도록 한다.

\section{수요와 공급 데이터 생성 원칙}

\subsection{명목시장과 실제시장 구분}

본 연구에서는 Zoe, Chiara(2018)이 사용한 것처럼 소비자가 주문을 게시한 뒤 판매자가 그 주문을 보고 즉시 공급에 참여하는
완전 반응형 P2P 플랫폼 구조를 사용하지 않는다.
그러한 구조는 실제 수요가 발생한 이후 공급자가 반응할 수 있기 때문에
실제 수요를 알기 전에 운송용량을 선제적으로 확보해야 하는 1단계 robust optimization의 필요성이 약해질 수 있다.

따라서 기준 시나리오는 다음의 시간적 순서를 따른다.

\begin{enumerate}[label=\arabic*)]
    \item 시장에는 판매자별 품목 공급가능량 $\hat s_{ilt}$가 존재한다.
    \item 플랫폼은 과거자료 또는 예측을 통해 소비자별 명목수요 $\hat d_{jlt}$를 알고 있다.
    \item 실제 수요폭증과 운송망 충격이 실현되기 전에 플랫폼이 선제계약량 $x^R_{iklt}$를 결정한다.
    \item 이후 실행 시나리오가 실현되면 실제 수요와 사용가능 운송용량을 바탕으로
    $y_{ijklts}$, $m_{jlts}$, $z_{ijklts}$(실제 운송 물량, 이월 수요, 거래 성립)가 결정된다.
\end{enumerate}

즉, 1단계의 명목수요는 ``이미 확정된 주문''이라기보다
\textbf{선제계약 의사결정을 위한 예측수요}의 의미를 가지고, 2단계의 불확실성을 모르는 상황이다.
2단계에서는 실제 불확실성을 가지는 상황이 관측된 뒤 이에 맞춰 매칭을 조정한다.

\subsection{수요ㆍ공급 범위의 설정}

판매자별 품목 공급량과 소비자별 품목 수요량은 일정한 범위 또는 확률분포 안에서 생성한다.
다만 해당 범위를 처음부터 임의로 고정하기보다,
먼저 $\Gamma=0$인 nominal baseline에서 시장이 정상적으로 작동하도록 calibration한 뒤 확정한다.

초기 calibration에서 확인해야 할 조건은 다음과 같다.

\begin{itemize}
    \item 판매자 총공급이 수요에 비해 지나치게 크지 않아야 한다.
    \item 반대로 총공급이 항상 수요보다 부족하여 구조적인 backlog가 발생해서도 안 된다.
    \item 중량 및 부피 운송용량이 항상 여유롭거나 항상 부족하지 않아야 한다.
    \item 경매주기 $\Delta$가 달라질 때 매칭성과가 실제로 달라질 수 있는 중간 수준의 시장을 구성해야 한다.
\end{itemize}

수요와 공급이 한쪽으로 지나치게 치우치는 것을 방지하기 위해,
품목별 총공급과 명목수요의 비율도 함께 확인한다. 품목 $l$의 기준시장에서

\[
\rho_l^{SD}
=
\frac{\sum_i \hat s_{ilt}}{\sum_j \hat d_{jlt}}
\]

를 계산하고, nominal baseline에서는 수요와 공급이 대체로 균형을 이루면서도
실행단계의 수요폭증을 일부 흡수할 수 있도록 \textbf{약 1.1--1.3 수준을 초기 calibration 범위}로 사용한다.
이 범위는 현실시장의 고정된 비율을 의미하는 것이 아니라,
과잉공급과 구조적 공급부족을 동시에 피하기 위한 수치실험상의 시작점이며,
pilot test에서 서비스율, backlog 및 제약의 binding 여부를 확인한 뒤 최종 범위를 조정한다.

따라서 데이터 생성의 목적은 특정 현실시장의 Load-to-Truck Ratio를 복제하는 것이 아니라,
\textbf{수급 균형이 매칭성과에 영향을 줄 수 있는 통제된 market}을 만드는 데 있다.

\section{시장 시나리오와 실행 시나리오의 구분}

본 연구에서는 ``시장 시나리오''와 2단계 실행 시나리오 $s$를 구분한다.

\begin{table}[h!]
\centering
\caption{두 종류의 시나리오 구분}
\begin{tabular}{p{0.25\textwidth}p{0.30\textwidth}p{0.35\textwidth}}
\toprule
\textbf{구분} & \textbf{의미} & \textbf{예시} \\
\midrule
시장 시나리오 & 실험 시작 전에 정하는 시장의 구조적 상태 & Supply-rich, Balanced, Tight, Stress \\
실행 시나리오 $s$ & 선제계약 후 실제 운영에서 발생하는 불확실성의 실현 & Normal, Demand Surge, Network Disruption \\
\bottomrule
\end{tabular}
\end{table}

시장 시나리오는 판매자, 소비자, 운송업체 수와 품목별 공급ㆍ수요ㆍ운송용량을 어떻게 구성했는가를 나타낸다.
반면 실행 시나리오 $s$는 동일한 시장에서 선제계약이 끝난 뒤
수요가 증가하거나 운송용량이 감소하는 상황을 나타낸다.

이 구분이 중요한 이유는 경매주기 $\Delta$를 결정하는 시점 때문이다.
플랫폼은 미래에 $s_1$, $s_2$, $s_3$ 중 어느 상황이 발생할지 미리 알 수 없으므로,
실행 시나리오마다 서로 다른 경매주기를 사전에 선택하는 것은 적절하지 않다.
대신 \textbf{하나의 시장 instance 안에서 가능한 실행 시나리오를 모두 고려하여
그 시장에 대한 하나의 최적 경매주기}를 선택한다.

\section{Robust Optimization과 $\Gamma$의 역할}

\subsection{왜 강건성은 1단계 선제계약에 적용하는가}

강건성이 1단계에 들어가는 이유는 의사결정 시점에 있다.
선제계약량 $x^R_{iklt}$는 실제 수요폭증과 운송망 충격을 관측하기 전에 결정된다.
따라서 플랫폼은 미래 수요가 정확히 얼마가 될지 알지 못한 상태에서
얼마만큼의 운송용량을 미리 확보할지 결정해야 한다.

이때 $\Gamma$는 명목수요만을 그대로 믿을 것인지,
아니면 일부 소비자의 수요가 동시에 불리한 방향으로 변하는 상황까지
추가로 방어할 것인지를 조절한다.
반면 2단계의 실제 매칭결정은 수요와 운송망 상태가 실현된 이후 이루어지므로,
이미 관측된 상황에 대응하는 recourse decision으로 처리한다.

따라서 전체 구조는 다음과 같이 해석할 수 있다.

\[
\boxed{
\text{Robust pre-contracting before uncertainty}
\;\longrightarrow\;
\text{Adaptive matching after realization}
}
\]

\subsection{$\Gamma$ 실험설계}

$\Gamma$의 기준값을 사전에 하나로 고정하여 ``정답''처럼 사용하지 않는다.
먼저

\[
\Gamma=0
\]

을 nominal 또는 deterministic baseline으로 사용하여
강건한 추가보호가 없는 상태에서 시장을 정상화한다.

이후 소비자 수를 $n=|J|$라 할 때, 특정 중간값 하나만 임의로 선택하기보다
$0\leq\Gamma\leq n$의 전체 보호수준을 균등하게 확인할 수 있도록

\[
\boxed{
\Gamma/n\in\{0,0.25,0.50,0.75,1.00\}
}
\]

의 다섯 수준을 비교한다.
기준시장에서는 $|J|=8$이므로

\[
\boxed{\Gamma\in\{0,2,4,6,8\}}
\]

이 된다. 여기서 25\%, 50\%, 75\%는 실제 충격의 발생확률을 의미하지 않는다.
이는 uncertainty budget의 전체 범위를 일정한 간격으로 탐색하여
보호수준 증가에 따른 선제계약량과 목적함수의 변화를 체계적으로 비교하기 위한 sensitivity grid이다.

\begin{table}[h!]
\centering
\caption{강건성 수준의 해석}
\begin{tabular}{ccc}
\toprule
\textbf{수준} & \textbf{$\Gamma$} & \textbf{해석} \\
\midrule
Nominal & $0$ & 추가적인 robust protection 없음 \\
Low protection & $0.25n$ & 낮은 수준의 동시 수요편차 보호 \\
Moderate protection & $0.50n$ & 중간 수준의 동시 수요편차 보호 \\
High protection & $0.75n$ & 높은 수준의 동시 수요편차 보호 \\
Full protection & $n$ & 모든 소비자의 불리한 수요편차까지 보호 \\
\bottomrule
\end{tabular}
\end{table}

본 실험에서는 $\Gamma$ 변화에 따라 최소한 다음 결과가 변하는지 확인한다.

\[
\sum_{i,k,l,t}x^R_{iklt}, \qquad
\text{Objective Value}, \qquad
\text{Service Rate}, \qquad
\text{Backlog}.
\]

특히 기존 실험에서 $\Gamma$가 결과에 영향을 주지 않는 현상이 관찰되었으므로,
본격적인 시나리오 확장 이전에 robust 제약의 작동 여부를 검증해야 한다.

\subsection{$\Gamma$가 영향을 주지 않을 때의 진단}

$\Gamma$가 달라져도 결과가 동일한 경우에는
``강건성이 필요 없다''고 바로 결론내리지 않는다.
먼저 robust target과 실제 선제계약량의 관계를 확인한다.

각 $(l,t)$에서 robust target을

\[
R_{lt}(\Gamma)
=
\sum_j \hat d_{jlt}
+
\Gamma\eta_{lt}
+
\sum_j\mu_{jlt}
\]

라고 하고, 실제 총 선제계약량을

\[
X^R_{lt}=\sum_{i,k}x^R_{iklt}
\]

라고 할 때,

\[
\mathrm{Slack}^{R}_{lt}
=
X^R_{lt}-R_{lt}(\Gamma)
\]

를 출력한다.

\begin{table}[h!]
\centering
\caption{$\Gamma$ 무반응 현상의 진단 기준}
\begin{tabular}{p{0.34\textwidth}p{0.50\textwidth}}
\toprule
\textbf{관찰결과} & \textbf{해석} \\
\midrule
$R_{lt}(\Gamma)$ 자체가 변하지 않음 & $\Gamma$ 전달 또는 수요편차 생성ㆍ집계 로직을 우선 점검 \\
$R_{lt}(\Gamma)$는 변하지만 $X^R_{lt}$가 동일 & robust 하한보다 이미 많은 선제계약을 하고 있어 robust 제약이 non-binding일 가능성 \\
$X^R_{lt}$도 변하지만 최종 목적값이 거의 동일 & 2단계의 공급 또는 운송용량이 충분하여 추가 선제계약의 경제적 효과가 작을 가능성 \\
\bottomrule
\end{tabular}
\end{table}

\section{선제계약비용과 과잉계약 문제}

강건성의 효과를 확인하려면 선제계약이 무조건 많이 이루어지는 구조를 피해야 한다.
선제계약비용 $h^R$이 지나치게 낮으면,
플랫폼은 $\Gamma$와 관계없이 가능한 한 많은 용량을 미리 계약하는 것이 유리할 수 있다.
이 경우 robust target이 증가해도 이미 충분히 높은 $x^R$를 선택하고 있으므로
$\Gamma$가 결과에 영향을 주지 않는다.

반대로 선제계약비용이 지나치게 높으면,
플랫폼은 미래 충격에 대비하기보다 미충족을 감수하는 쪽을 선택할 수 있다.
따라서 목표는 단순히 $h^R$을 ``충격 penalty보다 크게'' 또는 ``작게'' 만드는 것이 아니라,
\textbf{한 단위의 추가 선제계약이 가져오는 기대편익과 계약비용 사이에 실제 trade-off가 발생하도록}
비용수준을 calibration하는 것이다.

이를 위해 초기 실험에서는 seed별로 한 번 생성한 기준 선제계약비용 $h^{R,base}_{iklt}$를 고정하고,
비용배율 $\lambda_h$만 변경하는 sensitivity analysis를 수행한다.

\[
h^{R,test}_{iklt}=\lambda_h h^{R,base}_{iklt},
\qquad
\lambda_h\in\{0.5,1,2,4\}.
\]

초기에는 비교적 넓은 배율을 사용하여 해가 실제로 변하는 구간을 찾고,
필요할 경우 변화가 집중되는 구간을 더 세분화하여 추가 실험한다.
이는 새로운 모형구조를 추가하기 위한 것이 아니라,
어느 비용구간에서 $\Gamma$가 선제계약량을 실제로 변화시키는지 확인하고
선제계약비용과 robust protection 간 trade-off를 검증하기 위한
\textbf{모형 calibration 및 sensitivity analysis 절차}이다.

\section{서비스율 제약의 검토}

서비스율 제약은 목적함수만으로는 일부 수요가 과도하게 희생될 가능성을 막는
최소 운영기준의 역할을 한다.
다만 서비스율 하한이 지나치게 높으면
해당 제약 자체가 많은 선제계약과 운송을 강제하여,
robust 제약의 효과가 관찰되지 않을 수 있다.

따라서 서비스율 제약을 즉시 제거하기보다는 다음과 같이 단계적으로 검토한다.

\begin{enumerate}[label=\arabic*)]
    \item 기준시장 $\Gamma=0$에서 서비스율 제약의 slack을 확인한다.
    \item 서비스율 제약이 항상 binding인지 확인한다.
    \item 필요할 경우 $SR^{\min}$ 수준을 몇 개의 값으로 변화시켜 $\Gamma$와의 상호작용을 확인한다.
    \item 최종 연구에서는 서비스율을 핵심 목적함수가 아니라 최소 운영가능성 조건과 결과 KPI로 해석한다.
\end{enumerate}

이 과정의 목적은 서비스율을 연구의 중심으로 끌어오는 것이 아니라,
서비스율 제약이 선제계약 및 최적 경매주기 선택을 인위적으로 지배하고 있지 않은지 확인하는 데 있다.

\section{경매주기와 분석기간}

경매주기 $\Delta$는 MILP 내부에서 직접 최적화하는 의사결정변수가 아니라
outer-loop에서 후보를 하나씩 평가하는 설계변수이다.
각 $\Delta$에 대해 일별 원자료를 해당 경매슬롯 단위로 재집계한 뒤,
동일한 시장과 동일한 불확실성 조건에서 two-stage MILP를 해결한다.

각 시장 instance에 대한 최적 경매주기는

\[
\boxed{
\Delta^*
=
\arg\max_{\Delta\in\mathcal D}
\Phi(\Delta)
}
\]

로 정의한다.
여기서 $\Phi(\Delta)$는 해당 경매주기에서 모형이 산출한 운영성과 목적함수 값이다.

분석기간은 기존 30일 또는 60일보다 긴 기간을 사용하여
수요충격, backlog 누적, 이후 회복과정을 충분히 관찰하는 것이 적절하다.
최종 실험의 기준 calendar horizon으로

\[
\boxed{\mathcal T^{cal}=140\text{ days}}
\]

를 고려한다.
140일은 20주에 해당하므로 추후 주 단위 패턴을 도입할 경우에도 확장성이 있다.
다만 현재 기준 시나리오에서는 요일 선호나 지역효과를 사용하지 않으므로,
140일을 사용하는 주된 이유는 \textbf{장기간의 누적효과와 여러 경매주기의 반복성과를 안정적으로 관찰하기 위함}이다.

\section{Seed와 반복 시장의 역할}

Seed는 새로운 종류의 시장 시나리오가 아니라,
동일한 시장구조에서 난수로 생성되는 개별 수요ㆍ공급ㆍ비용 데이터를 재현하기 위한
\textbf{난수 생성 시작값}이다.

하나의 seed가 정해지면 해당 market instance의 일별 명목수요, 판매자 공급량,
운송업체의 중량ㆍ부피 capacity, 가치ㆍ비용계수, 수요편차 및 실행 시나리오의 난수 실현값 등
모든 외생적 raw data를 한 번 생성한 뒤 고정한다.
이후 $\Gamma$, $\Delta$ 또는 선제계약비용의 sensitivity 배율을 비교할 때는
비교대상이 되는 실험요인만 변경하고 나머지 raw data는 동일하게 유지한다.
특히 $\Delta$가 달라져도 새로운 데이터를 생성하지 않고,
동일한 일별 원자료를 각 경매주기에 맞게 재집계한다.

예를 들어 동일한 Balanced 시장이라도 seed가 달라지면
개별 소비자 수요, 판매자 공급량, 운송업체 capacity의 세부값은 달라질 수 있다.
그러나 Balanced라는 시장의 전체적 수급특성은 유지되어야 한다.
또한 동일 seed를 서로 다른 market regime에 적용할 경우에는 underlying random pattern은 유지하되,
해당 regime을 구성하기 위한 공급ㆍ수요ㆍcapacity의 구조적 scaling만 다르게 적용한다.

따라서 최종 단계에서 여러 시장상태를 비교할 때는
동일한 seed 집합을 각 시장상태에 반복 적용하는 방식이 적절하다.

예를 들어

\[
\mathcal Q=\{101,211,307,401,503\}
\]

를 사용할 경우,
Supply-rich, Balanced, Tight, Stress 각각에 동일한 5개 seed를 적용하여
총 20개의 synthetic market instance를 생성할 수 있다.

이 과정은 최초 baseline validation이 끝난 뒤 실시한다.
처음부터 20개 시장을 생성하면 $\Gamma$ 무반응 현상의 원인이
모형구조인지 데이터 생성인지 구분하기 어렵기 때문이다.

\section{전체 실험 순서}

본 연구의 실험은 한 번에 모든 시장조건을 변화시키는 방식보다
\textbf{모형 정상화 $\rightarrow$ 수급 calibration $\rightarrow$ 강건성 검증 $\rightarrow$
선제계약비용 sensitivity analysis $\rightarrow$ 경매주기 최적화 $\rightarrow$ 시장확장}
순으로 진행한다.

\subsection{1단계: Baseline Market Calibration}

첫 단계에서는 하나의 기준시장만 사용한다.
참여자 규모는 $|I|=6$, $|J|=8$, $|K|=4$, $|L|=2$로 고정하고,
판매자 3명은 품목 $A$, 나머지 3명은 품목 $B$를 공급하도록 한다.

이 단계에서는 $\Gamma=0$으로 설정하여
강건한 추가보호가 없는 nominal market을 먼저 정상화한다.
수요, 공급, 중량ㆍ부피 capacity 및 선제계약비용을 조정하면서
다음 조건을 만족하는 데이터를 찾는다.

\begin{itemize}
    \item 대부분의 경우 feasible solution이 존재한다.
    \item 모든 수요가 항상 100\% 충족되는 과잉공급 시장이 아니다.
    \item 반대로 모든 경매주기에서 구조적으로 큰 backlog가 발생하는 과소공급 시장도 아니다.
    \item 짧은 경매주기와 긴 경매주기 사이에 목적함수의 차이가 발생한다.
    \item 선제계약량이 계약상한에 항상 붙어 있지 않는다.
\end{itemize}

즉, 첫 단계의 목적은 ``최적 결과를 얻는 것''이 아니라
\textbf{이후 실험이 의미 있게 반응할 수 있는 중간 수준의 기준시장을 만드는 것}이다.

\subsection{2단계: $\Gamma$ Validation}

기준시장이 정상화되면 동일한 시장 데이터에서

\[
\Gamma=0,\quad 2,\quad 4,\quad 6,\quad 8
\]

을 적용한다.
이때 가장 먼저 확인할 것은 최종 목적함수 자체가 아니라
robust target, 실제 선제계약량, robust constraint slack이다.

$\Gamma$ 증가에 따라 robust target이 증가하는지,
그리고 그 증가가 실제 $x^R$에 반영되는지를 확인한다.
반영되지 않는다면 선제계약비용, 서비스율 제약, 계약상한,
공급 및 운송 capacity 가운데 어느 요소가 선제계약을 지배하고 있는지 분석한다.

이 단계를 통과해야 이후 ``강건성 수준에 따른 성과 차이''를 해석할 수 있다.

\subsection{3단계: Pre-contracting Cost Sensitivity Analysis}

$\Gamma$가 의도한 대로 작동하는 것이 확인되면,
동일한 seed와 동일한 market instance에서 기준 선제계약비용을 유지한 채
$\lambda_h\in\{0.5,1,2,4\}$의 배율만 변경하여 선제계약비용의 민감도를 확인한다.
이 단계에서는 선제계약비용이 너무 낮아 과잉계약을 유도하거나,
반대로 너무 높아 robust protection 자체를 무의미하게 만드는 구간이 있는지를 확인한다.
초기 넓은 범위의 실험에서 변화구간이 확인되면 그 구간만 세분화하여 추가 분석한다.

\subsection{4단계: Auction Cycle Optimization}

$\Gamma$가 의도한 대로 작동하는 것이 확인되면
각 강건성 수준에서 경매주기 후보를 완전열거한다.

\[
\Delta\in\mathcal D.
\]

각 $\Delta$마다 동일한 일별 원자료를 재집계하여 MILP를 풀고,
목적함수 값을 비교한다.
그 결과 가장 큰 목적함수를 가지는 경매주기를
해당 시장 instance와 강건성 수준에서의 최적 경매주기 $\Delta^*$로 정의한다.

이 단계의 핵심은 특정 실행 시나리오 $s$ 하나에 대해서만 $\Delta^*$를 찾는 것이 아니라,
동일한 시장에서 발생 가능한 실행 시나리오를 함께 고려하여
\textbf{사전에 적용할 하나의 경매 운영주기}를 도출하는 것이다.

\subsection{5단계: Market Regime Expansion}

기준시장과 $\Gamma$--$\Delta$ 구조가 정상적으로 작동한 이후에
시장 수급상태를 확장한다.
시장상태는 예를 들어 다음과 같이 구성할 수 있다.

\begin{table}[h!]
\centering
\caption{확장실험의 시장상태}
\begin{tabular}{cl}
\toprule
\textbf{시장상태} & \textbf{의도한 해석} \\
\midrule
Supply-rich & 공급 및 운송 capacity에 상대적 여유가 있는 시장 \\
Balanced & 수요와 공급ㆍ운송 capacity가 비교적 균형적인 시장 \\
Tight & 수요가 공급 및 운송 capacity에 근접한 시장 \\
Stress & 수요가 공급 또는 운송 capacity를 압박하는 시장 \\
\bottomrule
\end{tabular}
\end{table}

각 시장상태는 지역 또는 운송업체 선호를 통해 만드는 것이 아니라,
\textbf{품목별 수요량, 공급량 및 운송 capacity의 상대적 규모}를 조정하여 구성한다.

각 상태에 동일한 5개 seed를 적용하면

\[
4\text{ market regimes}\times5\text{ seeds}
=
20\text{ market instances}
\]

가 된다.
각 market instance에서 $\Gamma$와 $\Delta$를 동일한 방식으로 평가한다.

\subsection{6단계: Cross-Market Evaluation}

각 market instance에서 도출된 최적 경매주기를 비교하여
시장상태에 따라 $\Delta^*$가 체계적으로 달라지는지 확인한다.
또한 특정 시장에서 선택한 경매주기를 다른 seed 또는 다른 시장에 적용했을 때의
opportunity loss를 Regret으로 평가할 수 있다.

이 단계는 ``평균 목적함수가 가장 크다''는 사실만으로 강건성을 주장하기보다,
선택된 경매주기가 다른 synthetic market에서도 얼마나 안정적으로 성과를 유지하는지를 확인하기 위한 단계이다.

\section{최종 실험구조 요약}

\begin{longtable}{p{0.23\textwidth}p{0.62\textwidth}}
\caption{최종 실험설계 요약} \\
\toprule
\textbf{항목} & \textbf{설정 및 원칙} \\
\midrule
\endfirsthead
\toprule
\textbf{항목} & \textbf{설정 및 원칙} \\
\midrule
\endhead

연구초점 & 시장논리에 따른 N:M matching, robust 선제계약, 최적 경매주기 도출 \\
참여자 비율 & 판매자 : 소비자 : 운송업체 $=1.5:2:1$, 기준 규모 $6:8:4$ \\
품목 수 & 2개. 품목별 우선순위 없음 \\
품목 차이 & 중량ㆍ부피 소모특성만 차별화 \\
판매자--품목 관계 & 판매자별 전문품목을 부여하고 비취급 품목의 공급량을 0으로 설정 \\
지역ㆍ거리 & 모형에서 제외 \\
가격 & Core MILP의 주된 연구대상으로 사용하지 않고 MMTR--MVCG 사후 가격청산 유지 \\
1단계 & 불확실성 실현 이전 robust pre-contracting \\
2단계 & 실제 수요 및 운송망 상태 관측 이후 adaptive matching \\
강건성 수준 & $\Gamma/n=0,0.25,0.50,0.75,1.00$ 비교. 기준시장에서는 $0,2,4,6,8$ \\
경매주기 & 동일 일별 데이터에 대해 후보 $\Delta$를 outer-loop에서 비교 \\
Calendar horizon & 최종 실험에서 140일 고려 \\
Seed & 동일 seed에서는 raw data를 고정하고 실험요인만 변경하여 반복성과 재현성 확보 \\
시장 확장 & Baseline validation 이후 Supply-rich, Balanced, Tight, Stress로 확장 \\
핵심 결과 & 선제계약량, 목적함수, 서비스율, backlog, 최적 경매주기, Regret \\
\bottomrule
\end{longtable}

\section{연구 시나리오의 최종 해석}

본 연구의 시나리오는 ``현실의 특정 지역 화물시장''을 모사하는 것이 아니라,
5PL 플랫폼의 핵심 구조를 통제된 환경에서 검증하기 위한 synthetic market이다.
따라서 외부 지역선호, 노선선호, 품목우선순위와 같은 요소를 제거하고
판매자 공급, 소비자 수요, 운송업체 capacity라는 세 요소의 상대적 균형에 집중한다.

플랫폼은 실제 수요충격과 운송망 충격을 알기 전에
robust optimization을 이용하여 선제계약량을 정한다.
이후 실제 상황이 실현되면 N:M matching을 수행한다.
동일한 시장에 대해 경매주기 $\Delta$를 바꾸어 반복적으로 최적화하고,
목적함수를 최대화하는 $\Delta^*$를 해당 시장의 최적 경매주기로 해석한다.

실험은 처음부터 복잡한 시장시나리오를 다수 생성하는 방식으로 시작하지 않는다.
먼저 하나의 baseline market에서 $\Gamma=0$의 nominal case를 정상화하고,
수요와 공급이 지나치게 치우치지 않도록 수급수준을 calibration한다.
그 다음 $\Gamma/n=0,0.25,0.50,0.75,1.00$을 통해 robust 제약이 실제 선제계약에 영향을 주는지 검증하고,
선제계약비용에 대한 sensitivity analysis를 수행한다.
이후 경매주기 최적화를 수행하고,
마지막으로 서로 다른 수급상태와 seed로 시장을 확장한다.

따라서 본 연구의 실험 논리는 다음과 같이 요약할 수 있다.

\[
\boxed{
\text{Baseline Market Calibration}
\rightarrow
\Gamma\text{ Validation}
\rightarrow
h^R\text{ Sensitivity}
\rightarrow
\Delta\text{ Optimization}
\\
\rightarrow
\text{Market Regime Expansion}
\rightarrow
\text{Cross-Market Evaluation}
}
\]

이 순서를 따르면 특정 실험결과가 모형의 구조에서 발생한 것인지,
강건성 수준에서 발생한 것인지,
경매주기에서 발생한 것인지,
또는 단순한 random seed의 우연에서 발생한 것인지를 단계적으로 분리하여 해석할 수 있다.

\end{document}
